\section{はじめに}
\label{sec:introduction}

コート幾何は、テニス映像の下流処理すべての基準座標系である。
ボールの三角測量、選手のコート座標への写像、カメラ間の統合、
そして物理的に意味のある軌跡の推定は、
画像中の画素が規制コート上のどこへ対応するかを知っていることを前提にする。
しかし実際の映像では、カメラは移動し、コートを切り取り、斜めから見て、
ときには複数のコートを同時に映す。
したがって有用なコート検出器は、
見えている白線の断片に印を付けるだけでなく、
それらを一つの大域的説明へまとめる必要がある。

安価な注釈は、しばしばその逆を促す。
ヒートマップの教師は交点がどの画素にあるかを与えるが、
その交点がなぜその画素にあるかは与えない。
局所的に似た二つの線の交わりは、
別々の物理的な点、別々のコート、別々のカメラ姿勢に対応しうる。
大容量のデコーダは、データセットの局所的なテクスチャ統計を学習しながら、
下流が必要とする射影構造を獲得しないままでいられる。
画像ごとのカメラ校正を付ければこの曖昧性は消えるが、
学習用の全フレームにメートル系の外部・内部パラメータを付与する費用は、
少数の点をクリックする費用とは比較にならない。

\paragraph{着想}
本研究の出発点は、
再構成された三次元シーンを、合成らしい背景の供給源としてではなく
\emph{幾何教師の生成器}として使えるという観察である。
一つの動画から、リポジトリが公開する二つのコマンド境界を通して
3DGS シーン、校正済みの復元カメラ、微分可能なレンダラを得る\cite{kerbl2023gaussian}。
このシーンを規制コートへメートル単位で整合させたうえで、
制御された三次元軌道に新しいカメラを置けば、
各レンダ結果は厳密なカメラ行列、シーン--コート剛体変換、
各コート点の投影座標と描画支持情報を自動的に受け取る。
新しい視点を一つ増やす限界費用は、人手の注釈ではなく解析的な投影になる。

\paragraph{主たる問い}
本研究が評価したい問いは次の一つである。

\begin{quote}
実会場の外観を継承したまま、撮影位置・方向・可視範囲の分布を広げた新規視点を
学習へ加えることは、放送映像中心の実データでは不足しがちな視点条件に対して
テニスコート検出を適応させられるか。
また、その効果を画像数の増加と切り分けられるか。
\end{quote}

この問いに対して、本稿は次の三つを提示する。

\begin{enumerate}[leftmargin=1.6em]
  \item \textbf{視点分布を制御した再構成由来の教師生成。}
  復元カメラの水平凸包を解析的に内側へオフセットした領域に軌道全体を収める制約の下で、
  円・楕円・矩形・超楕円の軌道族を構成し、
  弧長一様にサンプルした視点を軌道グループ単位で分割する。
  生成候補の分布ではなく、採用後に学習へ渡る視点の分布を記述できる形にする。

  \item \textbf{単一のアライメントからの整合教師。}
  一つのシーン--コート変換から、キーポイント、領域分割、二値の線、
  線の意味分類、カメラ姿勢を同じ画像フレーム内で同時に生成する。
  この契約により、教師間でscope（対象コート・可視性・座標系）が食い違わない。

  \item \textbf{共通のアライメント評価手続き。}
  比較相手の検出結果にも同じ RANSAC ベースのアライメントを適用し、
  検出精度と幾何登録の精度を分けて報告する手続きを定義する。
  学習データやアーキテクチャが異なるモデル間で、
  幾何の良し悪しを一つの数値へ潰さない。
\end{enumerate}

\paragraph{実測した比較}
本稿はこの三つを設計で終わらせず、
公開実装ベースラインとの比較を実行して報告する。
domain ごとに決定論的 \resBenchDomainSamples 件（seed \resBenchSeed）を選び、
実写 held-out validation と合成 trajectory-group test の両方で、
生のキーポイントに同一の RANSAC アライメントを適用した。
推論は全 CPU で行った。
実写 validation では両者が近く、
\Method が PCK@0.01 と対応誤差の中央値で僅かに上回る。
合成 test では差が大きく、
ベースラインは有効な対応をほとんど返せない。
古典的なコートモデル当てはめも別標本の補助 probe として走らせ、
実写では動くが合成の multi-court 条件ではクラッシュと縮退が生じることを記録する。

\paragraph{主張しないこと}
本稿は、画像数を増やすこと、3DGS を使うこと、複数のヘッドを設けることを
それぞれ単独の新規性としては主張しない。
本稿が提示できる比較は、
アーキテクチャと学習データの両方が同時に異なる checkpoint と
公開実装ベースラインの比較である。
したがって合成側で大きな差が出ても、
それは 3DGS 利用の因果効果を分離しない。
差を「視点拡張が効いた」と読むことはしない。
実写のみで学習した同一アーキテクチャ、同一枚数の視点狭域・広域比較、
既知会場と未知会場の分離、複数 seed は追加実験の対象であり、
本稿はそれらを実施済みとして扱わない。
また、学習へ渡った視点分布そのものの集計図も本稿には含めない。
したがって結論は「汎化性能が高いことを証明した」ではなく、
\emph{合成分布への適応を実測し、
視点拡張仮説を切り分ける評価手続きを与えた}、である。

以降、\cref{sec:related_work} で関連研究を二本柱で整理し、
\cref{sec:problem_setting} で問題設定と座標規約、
\cref{sec:system} でシステム概要を述べる。
\cref{sec:alignment} が本稿の中心であり、
\cref{sec:view_generation} で視点生成と教師生成、
\cref{sec:detector} で検出器を簡潔に述べる。
\cref{sec:experiments} で比較条件、\cref{sec:results} で結果、
\cref{sec:discussion} で限界と倫理、\cref{sec:conclusion} で結論を述べる。
